\Needspace{18\baselineskip}
\section{Worked problems}

\Needspace{12\baselineskip}
\subsection{Union}
\textbf{Problem.} \emph{Union.} Derive Equation~\eqref{eq:ch43-union}. For $E=256,k=8$, solve the
smallest $B$ for which the expected union reaches at least 90\% of experts.

\textbf{Solution.} Set $E[1-(1-k/E)^B]\ge0.9E$, hence
$(1-k/E)^B\le0.1$ and
\[
B\ge\frac{\ln 0.1}{\ln(1-k/E)}.
\]
For $E=256,k=8$, this is $\ln(0.1)/\ln(31/32)\approx72.5$, so the smallest
integer batch is 73.

\Needspace{12\baselineskip}
\subsection{Occupancy}
\textbf{Problem.} \emph{Occupancy.} For $E=128,k=8,B=64$ and $U=126$, compute average rows per
active expert. Explain why this can still be a poor GEMM shape.

\textbf{Solution.} $Bk/U=64(8)/126\approx4.06$ rows per active expert. Four rows are far smaller than
the row dimension for which large accelerator GEMMs normally achieve high utilization;
padding and launch/dispatch overhead can dominate.

\Needspace{12\baselineskip}
\subsection{Parameters}
\textbf{Problem.} \emph{Parameters.} A layer has 64 experts of 20M parameters, top-4 routing and
100M shared parameters. Compute total and active parameters for one token.

\textbf{Solution.} Total: $100\text{M}+64(20\text{M})=1.38$B. Active per token:
$100\text{M}+4(20\text{M})=180$M.

\Needspace{12\baselineskip}
\subsection{Correctness}
\textbf{Problem.} \emph{Correctness.} Work the three-token routing trace in this chapter and show
the combine result if the grouped rows return in reverse expert order.

\textbf{Solution.} Physical return order is irrelevant if every row carries its token ID and expert ID.
The combine operation scatters each result to its original token and applies the gate
paired with that expert. A reverse return order therefore produces the same mathematical
sum; positional combine without IDs is incorrect.
Use the executable's scalar experts $f_e(x)=(e+1)(x+0.1)$ and tokens
$(1,2,-1)$. Token zero routes to $(e=2,g=0.7)$ and $(e=0,g=0.3)$:
$y_0=0.7(3.3)+0.3(1.1)=2.64$. Token one routes to $(1,0.6)$
and $(2,0.4)$: $y_1=0.6(4.2)+0.4(6.3)=5.04$. Token two
routes to $(2,0.8)$ and $(3,0.2)$:
$y_2=0.8(-2.7)+0.2(-3.6)=-2.88$.
Reversing expert-return order must preserve $(2.64,5.04,-2.88)$
within the declared floating-point tolerance. Sharing the scalar expert
function between the reference and grouped paths checks dispatch/combine,
not the correctness of an independent FFN kernel.

\Needspace{12\baselineskip}
\subsection{Distributed bytes}
\textbf{Problem.} \emph{Distributed bytes.} With $B=128,k=4,d=4096$ and BF16 activations, compute
the naive activation bytes in one dispatch direction.

\textbf{Solution.} $Bkd b_a=128(4)(4096)(2)=4{,}194{,}304$ bytes, exactly 4 MiB in one naive dispatch
direction, before protocol, alignment or duplicate-destination optimizations.

\Needspace{12\baselineskip}
\subsection{Capacity experiment}
\textbf{Problem.} \emph{Implementation.} Extend the lab with a per-expert capacity of two rows.
First implement dropping, then dropless overflow. Demonstrate why the dropping version
can make one token's output depend on its batch-mates.

\textbf{Solution.} A dropping implementation is batch-dependent: adding another token can fill an expert's
capacity before the original token's assignment is admitted. Dropless overflow preserves
the model semantics and converts overload into additional work/latency instead.
